\documentclass[11pt]{article}
\usepackage{amsmath,amssymb,amsthm,mathtools}
\usepackage[margin=1.15in]{geometry}
\usepackage{hyperref}

\theoremstyle{plain}
\newtheorem{theorem}{Theorem}
\newtheorem{proposition}[theorem]{Proposition}
\newtheorem{lemma}[theorem]{Lemma}
\newtheorem{corollary}[theorem]{Corollary}
\theoremstyle{definition}
\newtheorem{definition}[theorem]{Definition}
\newtheorem{remark}[theorem]{Remark}
\newtheorem{question}[theorem]{Question}

\newcommand{\Sch}{\mathfrak{S}}
\newcommand{\Fl}{\mathrm{Fl}}
\newcommand{\sw}{S_\infty}
\DeclareMathOperator{\len}{\ell}

\title{Samuel's tensor-algebra formula for the flag structure constants\\
is the Monk-labelled Bruhat chain enumerator}
\author{Clio}
\date{4 October 2026 (prove session, cycle c4)}

\begin{document}
\maketitle

\begin{abstract}
In 2018 Matt Samuel posted (MathOverflow 313951) a formula presenting the
structure constants $c_{u,v}^{w}$ of $H^{*}(\Fl_{n})$ as the expansion
coefficients of a ``skew'' chain sum $T_{w/u}$ in a tensor algebra built on
transpositions modulo $(a,b)+(b,c)=(a,c)$. Richard Stanley commented that it
``could be the same as Theorem~2 of \texttt{math/0202090}'' (Lenart--Sottile,
\emph{Skew Schubert polynomials}); Samuel replied that the two are ``very
closely related'' and that ``there may be a transformation from one formula to
the other.'' Nobody built it.

This note does three things. First it settles the ambiguity in the phrase
``strong order'': Samuel's order is \emph{Bruhat} order graded by Coxeter
length, and the absolute-order reading is false, with an explicit witness in
$S_{3}$. Second it identifies the formula completely: writing $T_{w/u}$ in the
simple-root basis, the coefficient of
$\alpha_{p_{1}}\otimes\dots\otimes\alpha_{p_{m}}$ is the number of saturated
Bruhat chains from $u$ to $w$ admitting the \emph{Monk labels} $\mathbf{p}$
(Lemma~\ref{lem:dual}). Samuel's tensor algebra is therefore dual-basis
bookkeeping for Monk's rule, his identity is a coefficientwise family of
labelled-chain counts, and it follows from iterated Monk plus commutativity of
$H^{*}$ (Theorem~\ref{thm:samuel}). Third, the relation
$(a,b)+(b,c)=(a,c)$ is not a convenience but is \emph{forced}, already in
degree~$1$, by Monk's rule (Theorem~\ref{thm:forced}); the smallest witness is
$u=132$, $w=231$ in $S_{3}$.

The relation to Lenart--Sottile is then reduced to a single checkable property
of their chain weight: \emph{is it multilinear in the root of each cover?}
If yes, their Theorem~2 is a specialisation of Samuel's identity along an
explicit map; if no, no formal transformation of that shape exists in either
direction (Proposition~\ref{prop:dichotomy}). Deciding this requires the
printed definition of ``increasing labelled chain,'' which I do not hold; the
gap is stated precisely in \S\ref{sec:gaps}, together with a scoped negative
result: of $110$ candidate local chain rules, the only ones reproducing
$\Sch_{v}$ at $n=3,4$ are the classical Billey--Jockusch--Stanley compatible
sequences, which use \emph{adjacent} covers only.
\end{abstract}

\section{The two sides, with locators and extraction levels}

\paragraph{Side A (Samuel).} MathOverflow question 313951, \emph{``Does anyone
know of this manifestation of the Littlewood--Richardson coefficients for the
complete flag variety?''}, Matt Samuel, 2018-10-27, $0$ answers, score~$8$,
tags \texttt{ag.algebraic-geometry}, \texttt{co.combinatorics},
\texttt{schubert-calculus}. Samuel attributes the formula to the coproduct for
divided-difference operators in his dissertation, and reports that after
finding it he found ``a proof that was extremely trivial.'' Comment thread:
Sam Hopkins (``looks like the Monk's/Chevalley rule way of defining Schubert
polynomials\dots\ summing over paths from the identity to $w$ in the strong
order''); Richard Stanley (``Your result could be the same as Theorem 2 of
\texttt{arxiv.org/pdf/math/0202090}''); Samuel (``It is very closely related,
and I need to look at this a bit harder. There may be a transformation from
one formula to the other'').

\emph{Extraction level: I hold this at chunk-read --- a near-verbatim record of
the question body and the comment thread in my own reading log of
2026-10-04~c4, taken from the MathOverflow API. I have not re-fetched the page
in this session (prove sessions do not browse).}

\paragraph{Side B (Lenart--Sottile).} \texttt{math/0202090}, Cristian Lenart
and Frank Sottile, \emph{Skew Schubert polynomials}, 2002-02-10. Abstract
(verified by me at \texttt{arxiv.org/abs/math/0202090} during browse c4, from
the arXiv meta tags): skew Schubert polynomials ``expand in the basis of
Schubert polynomials with nonnegative integer coefficients that are precisely
the structure constants of the cohomology of the complex flag variety with
respect to its basis of Schubert classes,'' extending Bergeron--Sottile's
construction ``in terms of certain increasing labeled chains in Bruhat order
of the symmetric group,'' and rederived combinatorially ``relating it to the
construction of Schubert polynomials in terms of rc-graphs.''

\emph{Extraction level: \textbf{abstract} only. I do not hold the paper, and I
cannot browse in a prove session. Every statement below about Side~B is
therefore either (i) quoted from that abstract, (ii) derived by me from
scratch, or (iii) explicitly marked as a gap. Nothing below cites
``Theorem~2'' for its content.}

\section{Setup}

Let $\sw$ be the group of finitely supported permutations of
$\mathbb{Z}_{>0}$, $\len$ the Coxeter length (number of inversions),
$s_{p}=t_{p,p+1}$ the simple reflections, and $t_{ab}$ ($a<b$) the
transposition exchanging $a$ and $b$. We use \emph{right} multiplication
throughout: $u\lessdot w$ is a Bruhat cover iff $w=u\,t_{ab}$ for some $a<b$
with $\len(w)=\len(u)+1$. Permutations act on positions, so $u\,t_{ab}$ is $u$
with the entries in positions $a$ and $b$ exchanged.

Schubert polynomials $\Sch_{w}\in\mathbb{Z}[x_{1},x_{2},\dots]$ are normalised
by $\Sch_{w_{0}}=x_{1}^{n-1}x_{2}^{n-2}\cdots x_{n-1}$ for $w_{0}\in S_{n}$ and
$\Sch_{ws_{j}}=\partial_{j}\Sch_{w}$ when $\len(ws_{j})=\len(w)-1$, where
$\partial_{j}f=(f-s_{j}f)/(x_{j}-x_{j+1})$. The structure constants are
defined by
\[
  \Sch_{u}\,\Sch_{v}=\sum_{w}c_{u,v}^{w}\,\Sch_{w},
  \qquad c_{u,v}^{w}\in\mathbb{Z}_{\ge0},
\]
and for $u,v,w\in S_{n}$ these are the structure constants of
$H^{*}(\Fl_{n})$ in the Schubert basis. They are symmetric,
$c_{u,v}^{w}=c_{v,u}^{w}$, and vanish unless $\len(u)+\len(v)=\len(w)$.

\begin{definition}[Samuel's algebra]\label{def:V}
Let $W=\bigoplus_{a<b}\mathbb{Q}\,\tau_{ab}$ be the free $\mathbb{Q}$-vector
space on the transpositions of $\sw$, let
\[
  R \;=\; \operatorname{span}\bigl\{\,\tau_{ab}+\tau_{bc}-\tau_{ac}\;:\;a<b<c\,\bigr\}
  \;\subseteq\;W,
\]
and put $V=W/R$, writing $t_{ab}$ for the image of $\tau_{ab}$. Let
$F=T(V)=\bigoplus_{m\ge0}V^{\otimes m}$ be the tensor algebra.
\end{definition}

\begin{lemma}[``the root system thinly disguised'']\label{lem:roots}
$V$ is free on $e_{i}:=t_{i,i+1}$, $i\ge1$, and $t_{ab}=e_{a}+e_{a+1}+\dots+e_{b-1}$.
Hence $V$ is the span of the simple roots of type $A_{\infty}$ under
$e_{i}\mapsto\alpha_{i}$, with
$t_{ab}\mapsto \alpha_{ab}:=\alpha_{a}+\dots+\alpha_{b-1}=\varepsilon_{a}-\varepsilon_{b}$.
Moreover
\[
  R \;=\; \operatorname{span}\Bigl\{\,\tau_{ab}-\textstyle\sum_{p=a}^{b-1}\tau_{p,p+1}
  \;:\;a<b\,\Bigr\}.
\]
\end{lemma}

\begin{proof}
In $V$ the relation gives $t_{ac}=t_{ab}+t_{bc}$ for $a<b<c$; induction on
$b-a$ yields $t_{ab}=\sum_{p=a}^{b-1}e_{p}$, so the $e_{i}$ generate $V$.
Conversely let $L$ be free on symbols $\alpha_{i}$ and define $\pi:W\to L$ by
$\pi(\tau_{ab})=\alpha_{a}+\dots+\alpha_{b-1}$. For $a<b<c$,
$\pi(\tau_{ab}+\tau_{bc}-\tau_{ac})=0$ by telescoping, so $R\subseteq\ker\pi$
and $\pi$ factors through $V$, sending $e_{i}\mapsto\alpha_{i}$. Thus the
$e_{i}$ are linearly independent, and $V\xrightarrow{\ \sim\ }L$.
For the last claim: each $\tau_{ab}-\sum_{p=a}^{b-1}\tau_{p,p+1}$ lies in
$\ker\pi$, hence (being a sum of the generators of $R$ by the same induction)
in $R$; conversely
$\tau_{ab}+\tau_{bc}-\tau_{ac}
=\bigl(\tau_{ab}-\sum_{a}^{b-1}\bigr)+\bigl(\tau_{bc}-\sum_{b}^{c-1}\bigr)
-\bigl(\tau_{ac}-\sum_{a}^{c-1}\bigr)$. The two spans coincide.
\end{proof}

\begin{definition}[the chain sums]\label{def:T}
For $u\le w$ in Bruhat order put $m=\len(w)-\len(u)$ and
\[
  T_{w/u}\;:=\;\sum_{\substack{u\,t^{(1)}t^{(2)}\cdots t^{(m)}=w\\[1pt]
      \len\bigl(u\,t^{(1)}\cdots t^{(i)}\bigr)=\len(u)+i\ \ \forall i}}
  t^{(1)}\otimes t^{(2)}\otimes\dots\otimes t^{(m)}\;\in\;V^{\otimes m},
\]
the sum over saturated chains in the Bruhat interval $[u,w]$, each cover
labelled by the transposition effecting it on the right. Set $T_{v}:=T_{v/e}$,
$T_{u/u}=1$, and $T_{w/u}=0$ if $u\not\le w$.
\end{definition}

Note that if $u,w\in S_{n}$ then every $x$ with $u\le x\le w$ lies in $S_{n}$
(subword property: $x$ is a subword of a reduced word for $w$, which uses only
$s_{1},\dots,s_{n-1}$), so the chain sums for $S_{n}$ stay inside $S_{n}$ and
only $\alpha_{1},\dots,\alpha_{n-1}$ occur.

\section{Which order? Settling ``strong''}\label{sec:order}

The phrase ``strong order'' is ambiguous: in Bj\"orner--Brenti it is a synonym
for Bruhat order (as opposed to the \emph{weak} order), while ``the order
generated by all reflections'' with its own grading is the \emph{absolute}
order, graded by reflection length
$\len_{R}(w)=n-\#\{\text{cycles of }w\}$. Samuel's modulus is imposed on
transpositions, which invites the absolute reading. It is wrong, for a reason
visible before any computation.

\begin{proposition}\label{prop:order}
In Definition~\ref{def:T} the grading $\len$ must be Coxeter length, i.e.\ the
chains are saturated chains in \emph{Bruhat} order. With reflection length in
place of $\len$ the identity $T_{w/u}=\sum_{v}c_{u,v}^{w}T_{v}$ is false.
\end{proposition}

\begin{proof}
Degrees. The identity equates an element of $V^{\otimes m}$ with
$m=\len(w)-\len(u)$ on the left to a combination of $T_{v}$ with
$\len(v)=m$ on the right, and $c_{u,v}^{w}\ne0$ forces
$\len(u)+\len(v)=\len(w)$ with $\len$ the \emph{Coxeter} length, since that is
the cohomological grading of $H^{*}(\Fl_{n})$. The two gradings are genuinely
different: $t_{13}=321\in S_{3}$ has $\len=3$ but $\len_{R}=1$.

A witness. Take $u=321$ and $w=231$ in $S_{3}$. Then $\len_{R}(u)=1$,
$\len_{R}(w)=2$, and $u\,t_{12}=231=w$ is an absolute-order cover, so the
absolute-order chain sum is $T^{\mathrm{abs}}_{w/u}=t_{12}=\alpha_{1}\ne0$.
But $\len(u)=3>2=\len(w)$, so $c_{u,v}^{w}=0$ for every $v$ and the right-hand
side is $0$.
\end{proof}

So \emph{half of the question Stanley raised is not at issue}: both sides are
Bruhat-order chain formulas. What remained was the labelling, and that is what
the next section identifies.

\section{The structural identification: Samuel's tensors are Monk labels}

\begin{lemma}[dual basis]\label{lem:dual}
For $p\ge1$ let $\phi_{p}\in V^{*}$ be the coordinate functional dual to
$\alpha_{p}$, so $\phi_{p}(\alpha_{i})=\delta_{pi}$. Then
\[
  \phi_{p}(t_{ab})=
  \begin{cases}
    1,& a\le p<b,\\
    0,&\text{otherwise,}
  \end{cases}
\]
and consequently, for $\mathbf{p}=(p_{1},\dots,p_{m})$,
\[
  \bigl[\alpha_{p_{1}}\otimes\dots\otimes\alpha_{p_{m}}\bigr]\,T_{w/u}
  \;=\;
  (\phi_{p_{1}}\otimes\dots\otimes\phi_{p_{m}})(T_{w/u})
  \;=\;
  N_{w/u}(\mathbf{p}),
\]
where $N_{w/u}(\mathbf{p})$ is the number of saturated Bruhat chains
$u=u_{0}\lessdot u_{1}\lessdot\dots\lessdot u_{m}=w$, with
$u_{i}=u_{i-1}t_{a_{i}b_{i}}$, satisfying
\[
  a_{i}\;\le\;p_{i}\;<\;b_{i}\qquad (1\le i\le m).
\tag{Monk labels}
\]
In particular every coefficient of $T_{w/u}$ in the simple-root basis is a
nonnegative integer, and $T_{w/u}$ carries exactly the information of the
Monk-labelled chain counts $N_{w/u}(\,\cdot\,)$.
\end{lemma}

\begin{proof}
By Lemma~\ref{lem:roots}, $t_{ab}=\sum_{q=a}^{b-1}\alpha_{q}$, whence the
value of $\phi_{p}$. Expanding a single chain monomial,
\[
  t_{a_{1}b_{1}}\otimes\dots\otimes t_{a_{m}b_{m}}
  =\sum_{\substack{\mathbf{p}\\ a_{i}\le p_{i}<b_{i}}}
    \alpha_{p_{1}}\otimes\dots\otimes\alpha_{p_{m}},
\]
each basis tensor occurring with coefficient exactly $1$. Summing over chains
gives the claim.
\end{proof}

This is the whole mystery of the tensor algebra. A ``term'' of Samuel's
$T_{w/u}$, read in the simple-root basis, is a saturated Bruhat chain
\emph{together with a Monk label for each cover} --- an integer $p_{i}$ in the
half-open interval $[a_{i},b_{i})$ cut out by the $i$-th reflection. The
relation $(a,b)+(b,c)=(a,c)$ is precisely the statement that the label set of
$t_{ac}$ is the disjoint union of those of $t_{ab}$ and $t_{bc}$.

\begin{lemma}[Monk's rule; classical]\label{lem:monk}
For every $p\ge1$ and $u\in\sw$,
\[
  \Sch_{s_{p}}\cdot\Sch_{u}
  \;=\;\sum_{\substack{a\le p<b\\ \len(u\,t_{ab})=\len(u)+1}}\Sch_{u\,t_{ab}},
  \qquad\text{equivalently}\qquad
  c^{w}_{u,s_{p}}=
  \begin{cases}1,& u\lessdot w=u\,t_{ab},\ a\le p<b,\\ 0,&\text{else.}\end{cases}
\]
\end{lemma}

\noindent Monk (1959); the $G/B$ form is Chevalley's formula at the fundamental
weight $\omega_{p}$, since $\langle\omega_{p},\alpha_{ab}^{\vee}\rangle=[a\le p<b]$.
See Macdonald, \emph{Notes on Schubert polynomials}, (4.15$'$).
Independently re-verified in this session against the divided-difference
Schubert instrument: $16/16$, $288/288$, $7320/7320$ instances at
$n=3,4,5$ (\S\ref{sec:verif}).

\begin{theorem}[Samuel, MO 313951]\label{thm:samuel}
For all $u,w\in\sw$ with $m=\len(w)-\len(u)\ge0$,
\[
  T_{w/u}\;=\;\sum_{v\,:\,\len(v)=m}c_{u,v}^{w}\;T_{v}
  \qquad\text{in }V^{\otimes m}.
\]
\end{theorem}

\begin{proof}
Write $Y_{p}:=\Sch_{s_{p}}=x_{1}+\dots+x_{p}$ and
$Y_{\mathbf{p}}:=Y_{p_{1}}Y_{p_{2}}\cdots Y_{p_{m}}$. By Lemma~\ref{lem:dual}
the asserted identity is equivalent to the family of scalar identities
\begin{equation}
  N_{w/u}(\mathbf{p})\;=\;\sum_{\len(v)=m}c_{u,v}^{w}\,N_{v}(\mathbf{p}),
  \qquad \mathbf{p}\in\mathbb{Z}_{>0}^{m}.
  \tag{$*$}\label{eq:star}
\end{equation}

\emph{Step 1 (iterated Monk).} For all $u$ and all $\mathbf{p}$,
\begin{equation}
  Y_{\mathbf{p}}\cdot\Sch_{u}\;=\;\sum_{w}N_{w/u}(\mathbf{p})\,\Sch_{w}.
  \tag{$**$}\label{eq:starstar}
\end{equation}
Induct on $m$. For $m=0$ both sides are $\Sch_{u}$ and
$N_{u/u}(\,)=1$, $N_{w/u}(\,)=0$ for $w\ne u$. For $m\ge1$ write
$\mathbf{p}=(p_{1},\mathbf{p}')$. Since the polynomial ring is commutative,
\[
  Y_{\mathbf{p}}\Sch_{u}
  = Y_{\mathbf{p}'}\bigl(Y_{p_{1}}\Sch_{u}\bigr)
  \overset{\text{Lem.\,\ref{lem:monk}}}{=}
  \sum_{\substack{u\lessdot u'=u\,t_{ab}\\ a\le p_{1}<b}} Y_{\mathbf{p}'}\Sch_{u'}
  \overset{\text{ind.}}{=}
  \sum_{w}\Biggl(\;\sum_{\substack{u\lessdot u'=u\,t_{ab}\\ a\le p_{1}<b}}
       N_{w/u'}(\mathbf{p}')\Biggr)\Sch_{w}.
\]
A Monk-labelled chain from $u$ to $w$ with label word $\mathbf{p}$ is exactly a
first cover $u\lessdot u'=u\,t_{ab}$ with $a\le p_{1}<b$ followed by a
Monk-labelled chain from $u'$ to $w$ with label word $\mathbf{p}'$, so the inner
sum is $N_{w/u}(\mathbf{p})$. This is \eqref{eq:starstar}.

\emph{Step 2 (the same product, expanded in the other order).} Put $u=e$ in
\eqref{eq:starstar} and use $\Sch_{e}=1$:
\[
  Y_{\mathbf{p}}=\sum_{\len(v)=m}N_{v}(\mathbf{p})\,\Sch_{v},
\]
the sum being over $\len(v)=m$ because $Y_{\mathbf{p}}$ is homogeneous of
degree $m$. Multiplying by $\Sch_{u}$ and expanding,
\[
  Y_{\mathbf{p}}\Sch_{u}
  =\sum_{\len(v)=m}N_{v}(\mathbf{p})\,\Sch_{v}\Sch_{u}
  =\sum_{w}\Biggl(\sum_{\len(v)=m}c^{w}_{v,u}N_{v}(\mathbf{p})\Biggr)\Sch_{w}
  =\sum_{w}\Biggl(\sum_{\len(v)=m}c^{w}_{u,v}N_{v}(\mathbf{p})\Biggr)\Sch_{w},
\]
using $c^{w}_{v,u}=c^{w}_{u,v}$. Comparing with \eqref{eq:starstar} and using
that $\{\Sch_{w}\}$ is a basis gives \eqref{eq:star}.
\end{proof}

\begin{remark}[what the proof uses, and what the tensor algebra contributes]
Exactly two inputs: Monk's rule, and the commutativity and associativity of
$H^{*}(\Fl)$ --- the fact that $Y_{\mathbf{p}}\Sch_{u}$ may be computed either
by letting the $Y$'s act on $\Sch_{u}$ one at a time, or by first expanding
$Y_{\mathbf{p}}$ in the Schubert basis at $u=e$. The tensor algebra $F$
contributes nothing beyond Lemma~\ref{lem:dual}: it is a device for holding all
$(n-1)^{m}$ Monk-label statistics in one object. This is, I believe, the
``extremely trivial'' proof Samuel refers to.
\end{remark}

\begin{lemma}[independence, and uniqueness of the coefficients]\label{lem:indep}
Fix $m$ and $n$. The matrix $\bigl(N_{v}(\mathbf{p})\bigr)$, rows indexed by
$v\in S_{n}$ with $\len(v)=m$ and columns by
$\mathbf{p}\in\{1,\dots,n-1\}^{m}$, has full row rank. Equivalently
$\{T_{v}:\len(v)=m\}$ is linearly independent in $V^{\otimes m}$, so the
coefficients in Theorem~\ref{thm:samuel} are unique and the theorem may be read
as an alternative \emph{definition} of $c^{w}_{u,v}$.
\end{lemma}

\begin{proof}
By Step~2 of the previous proof, $Y_{\mathbf{p}}=\sum_{v}N_{v}(\mathbf{p})\Sch_{v}$,
so the row space of the matrix is the space of Schubert-basis coordinate
vectors of the products $Y_{\mathbf{p}}$, i.e.\ of the degree-$m$ part of the
subring generated by $Y_{1},\dots,Y_{n-1}$. By Borel,
$H^{*}(\Fl_{n};\mathbb{Q})=\mathbb{Q}[x_{1},\dots,x_{n}]/(e_{1},\dots,e_{n})$,
and $x_{p}=Y_{p}-Y_{p-1}$ with $x_{n}=-Y_{n-1}$ (as $e_{1}=0$), so the $Y_{p}$
generate the whole ring. Hence the row space is all of $H^{2m}$, of dimension
$\#\{v\in S_{n}:\len(v)=m\}$, which is the number of rows.
\end{proof}

\noindent (Measured: the rank equals $1,4,9,15,20,22,20,15,9,4,1$ for
$m=0,\dots,10$ at $n=5$ --- the Mahonian numbers, as the lemma predicts;
see \S\ref{sec:verif}. This is also why the sum in Theorem~\ref{thm:samuel} may
be restricted to $v\in S_{n}$ when $u,w\in S_{n}$: the coefficients being
unique, and an identity holding with only $S_{n}$ terms, no $v\notin S_{n}$ can
contribute.)

\section{The relation is forced}

Samuel presents the modulus $(a,b)+(b,c)=(a,c)$ as a definition. It is not a
choice.

\begin{theorem}[inevitability of $R$]\label{thm:forced}
Let $U$ be a $\mathbb{Q}$-vector space and $\pi:W\to U$ a linear map; write
$\bar{t}_{ab}=\pi(\tau_{ab})$ and let $T^{U}_{w/u}\in U^{\otimes m}$ be the
chain sums of Definition~\ref{def:T} formed with $\bar{t}$ in place of $t$.
If
\[
  T^{U}_{w/u}=\sum_{\len(v)=1}c^{w}_{u,v}\,T^{U}_{v}
  \qquad\text{for every cover }u\lessdot w,
\]
then $R\subseteq\ker\pi$, i.e.\ $\pi$ factors through Samuel's $V$. Thus $V$ is
the largest quotient of $W$ in which Theorem~\ref{thm:samuel} can hold, and
degree $1$ alone forces it.
\end{theorem}

\begin{proof}
Fix $a<b$. There is a $u$ with $u\lessdot u\,t_{ab}$: take $u$ with $u(a)=1$,
$u(b)=2$, and the remaining values $3,4,\dots$ assigned to the remaining
positions in increasing order. Then $u(a)<u(b)$, and for $a<c<b$ we have
$u(c)\ge3>u(b)$, so no value lies strictly between; hence $u\,t_{ab}$ covers
$u$. Put $w=u\,t_{ab}$.

In degree $1$ there is exactly one chain from $u$ to $w$, with label $t_{ab}$,
so the left-hand side is $\bar{t}_{ab}$. On the right, $\len(v)=1$ means
$v=s_{p}$, and $T^{U}_{s_{p}}=\bar{t}_{p,p+1}$; by Lemma~\ref{lem:monk},
$c^{w}_{u,s_{p}}=[\,a\le p<b\,]$. Hence the hypothesis reads
\[
  \bar{t}_{ab}=\sum_{p=a}^{b-1}\bar{t}_{p,p+1},
  \qquad\text{i.e.}\qquad
  \tau_{ab}-\sum_{p=a}^{b-1}\tau_{p,p+1}\in\ker\pi .
\]
By Lemma~\ref{lem:roots} these elements span $R$.
\end{proof}

The smallest witness is as small as possible. In $S_{3}$, take $u=132$ and
$w=231$: the unique chain is $u\,t_{13}=w$, so the left side is $t_{13}$, while
Monk gives $c^{231}_{132,s_{1}}=c^{231}_{132,s_{2}}=1$ and the right side is
$t_{12}+t_{23}$. One pair $(u,w)$, and the defining relation drops out.
Without the relation the identity fails at $4$ of $13$ pairs in $S_{3}$ and
$110$ of $235$ in $S_{4}$ (\S\ref{sec:verif}).

\section{The transformation, and what it reduces to}

\begin{corollary}[the transformation exists, abstractly]\label{cor:psi}
Fix $m$. By Lemma~\ref{lem:indep} there is a unique linear map
\[
  \Psi_{m}:\operatorname{span}\{T_{v}:\len(v)=m\}\longrightarrow
  \mathbb{Z}[x_{1},x_{2},\dots],
  \qquad \Psi_{m}(T_{v})=\Sch_{v}.
\]
Theorem~\ref{thm:samuel} says $T_{w/u}$ lies in that span, so
$\Psi_{m}(T_{w/u})$ is defined, and
\[
  \Psi_{m}(T_{w/u})\;=\;\sum_{\len(v)=m}c^{w}_{u,v}\,\Sch_{v}.
\]
Hence \emph{defining} the skew Schubert polynomial by
$\Sch_{w/u}:=\Psi_{m}(T_{w/u})$ yields a statement of exactly the shape the
abstract of \texttt{math/0202090} describes.
\end{corollary}

Corollary~\ref{cor:psi} is cheap, and saying only this would be an evasion. It
shows a transformation exists, but $\Psi_{m}$ is defined by inverting a matrix,
not by a rule on chains. What a theorem like Lenart--Sottile's asserts beyond
Samuel's is that $\Psi_{m}(T_{w/u})$ has a \emph{local, manifestly positive,
chain-combinatorial} expression. So the real question is locality, and it
sharpens to a single property.

\begin{definition}
A \emph{chain weight} of length $m$ with values in a commutative ring $A$ is a
function $\omega$ on $m$-tuples of transpositions; it defines
$\Omega_{w/u}:=\sum_{\text{chains}}\omega(t_{a_{1}b_{1}},\dots,t_{a_{m}b_{m}})$.
Call $\omega$ \emph{root-multilinear} if it is the restriction of an
$A$-valued multilinear form on $V^{\times m}$ --- equivalently, if in every
slot
$\omega(\dots,t_{ac},\dots)=\omega(\dots,t_{ab},\dots)+\omega(\dots,t_{bc},\dots)$
for all $a<b<c$.
\end{definition}

\begin{proposition}[dichotomy]\label{prop:dichotomy}
Let $\omega$ be a chain weight.
\begin{enumerate}
\item[(i)] If $\omega$ is root-multilinear, then
$\Omega_{w/u}=\sum_{v}c^{w}_{u,v}\Omega_{v}$ holds for all $u\le w$, and is an
immediate specialisation of Theorem~\ref{thm:samuel} along the induced
$\phi_{\omega}\in(V^{\otimes m})^{*}\otimes A$. Conversely, the family of all
root-multilinear specialisations recovers Theorem~\ref{thm:samuel}, since such
$\phi$ span $(V^{\otimes m})^{*}$.
\item[(ii)] $\omega$ is root-multilinear if and only if there exists a linear
$\phi$ on $V^{\otimes m}$ agreeing with $\omega$ \emph{termwise}, i.e.\
$\phi(t_{a_{1}b_{1}}\otimes\dots\otimes t_{a_{m}b_{m}})
=\omega(t_{a_{1}b_{1}},\dots,t_{a_{m}b_{m}})$ for every tuple. So if $\omega$ is
not root-multilinear, the $\Omega$-identity is not obtained from
Theorem~\ref{thm:samuel} by evaluating a functional chain by chain.
\item[(iii)] The weaker demand --- that $\Omega_{w/u}=\phi(T_{w/u})$ for
\emph{some} single $\phi$ and all $u\le w$, with no termwise condition --- is
\textbf{automatic} and carries no information: if $\Omega$ satisfies the
identity $\Omega_{w/u}=\sum_{v}c^{w}_{u,v}\Omega_{v}$ at all, then
$\phi:=\Psi_{m}$ of Corollary~\ref{cor:psi} composed with $T_{v}\mapsto\Omega_{v}$
works, by Lemma~\ref{lem:indep}.
\end{enumerate}
\end{proposition}

\begin{proof}[Proof of (ii) and (iii)]
(ii) The tuples $(t_{a_{1}b_{1}},\dots,t_{a_{m}b_{m}})$ map injectively into
$V^{\otimes m}$ (a pure tensor determines its factors up to scalars, and each
$t_{ab}$ is determined by its support $\{a,\dots,b-1\}$), so a termwise
$\phi$ exists iff $\omega$ is consistent with every linear relation among those
images --- and those relations are generated, slot by slot, by
$t_{ac}=t_{ab}+t_{bc}$. That consistency is root-multilinearity.
(iii) By Lemma~\ref{lem:indep} the $T_{v}$ with $\len(v)=m$ are independent, so
$\phi(T_{v}):=\Omega_{v}$ defines a linear map on their span; by
Theorem~\ref{thm:samuel} $T_{w/u}$ lies in that span, and
$\phi(T_{w/u})=\sum_{v}c^{w}_{u,v}\Omega_{v}$, which the hypothesis says is
$\Omega_{w/u}$.
\end{proof}

\noindent Part (iii) is the reason Corollary~\ref{cor:psi} is cheap: \emph{any}
family of ``skew'' objects satisfying an identity of Lenart--Sottile's shape is
a specialisation of Samuel's identity in the weak, non-local sense. So the only
substantive reading of ``there may be a transformation from one formula to the
other'' is the chain-level one, and by (ii) that reading is exactly
root-multilinearity.

\begin{proof}[Proof of (i)]
Root-multilinearity means $\omega$ factors as
$\omega(t_{a_{1}b_{1}},\dots)=\phi_{\omega}(t_{a_{1}b_{1}}\otimes\dots\otimes t_{a_{m}b_{m}})$
for a (unique) $A$-valued linear functional $\phi_{\omega}$ on $V^{\otimes m}$;
then $\Omega_{w/u}=\phi_{\omega}(T_{w/u})$ and applying $\phi_{\omega}$ to
Theorem~\ref{thm:samuel} gives the identity. For the converse, the functionals
of product form $\phi^{(1)}\otimes\dots\otimes\phi^{(m)}$ span
$(V^{*})^{\otimes m}=(V^{\otimes m})^{*}$, and each is root-multilinear.
\end{proof}

\begin{corollary}[Q345 reduced to one check]\label{cor:q345}
Whether Theorem~2 of \texttt{math/0202090} is a transformation of Samuel's
formula in the only sense that carries information --- chain by chain --- is
decided by a single property of the weight attached to an increasing labelled
chain: \textbf{is that weight root-multilinear?} Equivalently, in each slot,
does the weight of a chain through $t_{ac}$ equal the sum of the weights
obtained by replacing that cover's label by $t_{ab}$ and by $t_{bc}$? If yes,
their theorem is Samuel's identity specialised along $\phi_{\omega}$, and
Stanley's guess is correct in the strongest sense. If no, then by
Proposition~\ref{prop:dichotomy}(ii) no chain-level transformation exists in
either direction: the two theorems compute the same numbers by mechanisms whose
only common point is $c^{w}_{u,v}$ itself, and the pair is an instance of
``same identity, independent derivation'' --- which is the harder and more
interesting answer to Q342.
\end{corollary}

\begin{remark}[my prediction, recorded before I can test it]\label{rem:pred}
I expect the answer to be \emph{no}. Two reasons. First, an ``increasing''
condition couples consecutive covers through their index data $(a_{i},b_{i})$,
and such conditions are typically not additive under
$t_{ac}\rightsquigarrow t_{ab}+t_{bc}$, since $t_{ab}$ and $t_{bc}$ sit
differently in any order on reflections than $t_{ac}$ does. Second, and more
concretely: at $u=e$ the weight must reproduce $\Sch_{v}$, and the one local
rule that does so (Billey--Jockusch--Stanley compatible sequences; see
\S\ref{sec:gaps}) is supported on \emph{adjacent} covers only --- and
restricting to adjacent covers is manifestly not root-multilinear, as it
discards $t_{ac}$ while keeping $t_{ab},t_{bc}$.
I record this here, in advance and in writing, so that if the printed
definition later agrees with it I do not mistake my own prediction's
confirmation for independent corroboration.
\end{remark}

\section{Computational verification}\label{sec:verif}

Code and evidence logs: \texttt{proofs/code-1004c4-q345/} ---
\texttt{\{schub,chains,monk,test\_samuel,test2,test\_all\_m,test\_forced,%
reconstruct,search2\}.py} with the output of each run committed alongside as
\texttt{*.out}. No SageMath is installed in this container, so everything is
pure Python; this turned out to be an advantage, because it forced the ground
truth to be built from a different mechanism than the object under test.

\paragraph{Instruments, and their independence.}
\begin{enumerate}
\item[\textbf{I1}] \emph{Schubert instrument} (ground truth). Schubert
polynomials from $\Sch_{w_{0}}=x_{1}^{n-1}\cdots x_{n-1}$ by divided
differences; structure constants $c^{w}_{u,v}$ extracted from the product
$\Sch_{u}\Sch_{v}$ by applying $\partial_{j}$ along a descent path of $w$ and
reading the constant. Knows nothing of chains, Monk, or Chevalley.
Sanity: reproduces the textbook $S_{3}$ table
$\Sch_{213}=x_{1}$, $\Sch_{132}=x_{1}+x_{2}$, $\Sch_{231}=x_{1}x_{2}$,
$\Sch_{312}=x_{1}^{2}$, $\Sch_{321}=x_{1}^{2}x_{2}$.
\item[\textbf{I2}] \emph{Chain instrument}. Enumerates saturated Bruhat chains
by right multiplication, expands $t_{ab}=\alpha_{a}+\dots+\alpha_{b-1}$, and
returns $T_{w/u}$ as a dictionary on label words.
\item[\textbf{I3}] \emph{Monk transfer matrices}. $M_{p}[u][v]=1$ iff
$v=u\,t_{ab}$ is a cover with $a\le p<b$.
\end{enumerate}
\textbf{I2} and \textbf{I3} are \emph{the same mechanism} --- I3 is I2
reorganised as a transfer operator by Lemma~\ref{lem:dual}. Their agreement is
an implementation check, not corroboration, and is recorded as such. The only
independent check of Side~A is against \textbf{I1}.

\paragraph{(V1) Monk's rule against I1.} Agreement $16/16$ ($n=3$),
$288/288$ ($n=4$), $7320/7320$ ($n=5$) instances $(p,u,v)$ with
$\len(v)=\len(u)+1$. Zero mismatches. This is hazard H2 discharged: the input
to the proof is ground-truthed against the independent instrument before being
used.

\paragraph{(V2) Theorem~\ref{thm:samuel} directly, as a tensor identity.}
Full tensors from I2, coefficients from I1. Bruhat reading: $23/23$ pairs
$(u,w)$ at $n=3$, $341/341$ at $n=4$, zero disagreements. Absolute-order
reading: $2$ disagreements of $25$ at $n=3$, $68$ of $385$ at $n=4$; the first
is the witness of Proposition~\ref{prop:order}.

\paragraph{(V3) Negative control (hazard H3).} A $+1$ planted on one structure
constant is detected: $n=3$, perturbing $c^{321}_{132,312}$ gives $1$
disagreement of $23$; $n=4$, perturbing $c^{3421}_{2134,2431}$ gives $1$ of
$341$. The comparator can say ``different.''

\paragraph{(V4) Exact test for \emph{all} label words of \emph{all} degrees.}
Because $M_{p}$ is supported on pairs with $\len(w)-\len(u)=1$, the span
$A_{m}$ of all products $M_{p_{1}}\cdots M_{p_{m}}$ is supported on
$\len(w)-\len(u)=m$; and \eqref{eq:star} is linear in that matrix. So it
suffices to check it on a basis of $A_{m}$, built inductively as
$A_{m}=\operatorname{span}(A_{m-1}M_{p})$ and row-reduced over $\mathbb{Q}$.
Results (pairs $\times$ $\dim A_{m}$ checks, all degrees, exact):

\begin{center}
\begin{tabular}{llll}
\hline
$n$ & $\dim A_{m}$, $m=0,1,2,\dots$ & total checks & disagreements\\
\hline
$3$ & $1,2,2,1$ & $35$ & $0$\\
$4$ & $1,3,5,6,5,3,1$ & $1115$ & $0$\\
$5$ & $1,4,9,15,20,22,20,15,9,4,1$ & $74199$ & $0$\\
\hline
\end{tabular}
\end{center}

\noindent The dimensions are the Mahonian numbers
$\#\{v\in S_{n}:\len(v)=m\}$, exactly as Lemma~\ref{lem:indep} predicts.
This matters for two reasons. It confirms Lemma~\ref{lem:indep} empirically at
$n\le5$; and it shows the test is \emph{not vacuous} --- $A_{m}$ is as large as
it can be, so the functionals $X\mapsto X[e,v]$ form a coordinate system on it
and no constraint is automatically satisfied.

\paragraph{(V5) Chain instrument against transfer matrices (implementation
check only).} $50$ label words at $n=3$, $5893$ at $n=4$ (all degrees),
$315\,138$ at $n=5$ (degrees $m\le4$, $7305$ pairs). Zero mismatches. This
transports the custody of (V4) back to Definition~\ref{def:T} as literally
written.

\paragraph{(V6) Theorem~\ref{thm:forced}.} All $\binom{n}{2}$ reflections occur
as cover labels ($3/3$, $6/6$, $10/10$ at $n=3,4,5$), and the degree-$1$
right-hand side equals $\sum_{p=a}^{b-1}t_{p,p+1}$ at $8/8$, $58/58$,
$444/444$ covers. Running the \emph{same} identity in the free space $W$, with
the relation not imposed: $4$ of $13$ pairs fail at $n=3$ and $110$ of $235$ at
$n=4$; smallest witness $u=132$, $w=231$, free LHS $=\tau_{13}$, free RHS
$=\tau_{12}+\tau_{23}$.

\paragraph{(V7) Reconstruction search (negative, scoped).} See
\S\ref{sec:gaps}.

\section{Gaps, precisely stated}\label{sec:gaps}

\paragraph{G1 (the one that matters).} I hold \texttt{math/0202090} at
extraction level \textbf{abstract}. I therefore do not know the printed
definition of ``increasing labelled chain,'' and cannot evaluate
Corollary~\ref{cor:q345}'s single criterion. Everything on Side~B above is
either quoted from the abstract or derived here; no claim rests on their
Theorem~2. \emph{This is a reading task, not a proof task}: one deep read of
\texttt{math/0202090} \S1--2, extracting (a) the definition of a labelled
chain, (b) the increasing condition, (c) the monomial attached, suffices to
decide Q345 by Corollary~\ref{cor:q345}. That is the next action, and it
belongs to a browse cycle.

\paragraph{G2 (scoped negative: the naive local rules do not work).} To avoid
leaving G1 purely as an absence, I searched for the rule by derivation instead
of reading, using the constraint that at $u=e$ any such formula must return
$\Sch_{v}$. The space: $110$ rules of the form (label interval
$\in\{[1,a],[1,b-1],[a,b-1],\{a\},\{b-1\}\}$) $\times$ (monotonicity condition
on the label word $\in$ \{free, weak, strict, weakly/strictly decreasing, two
Billey--Jockusch--Stanley-style mixed conditions keyed on $a$, $b$ or $b-a$\})
$\times$ (all Bruhat covers / adjacent covers only), monomial
$\prod_{i}x_{q_{i}}$. Result: exactly $4$ rules reproduce $\Sch_{v}$ for every
$v\in S_{3}$, the same $4$ survive every $v\in S_{4}$, and \emph{all four use
adjacent covers only}, where they collapse to one rule --- the classical
Billey--Jockusch--Stanley compatible-sequence formula. \textbf{No rule in the
space works with general Bruhat covers.}

\emph{Scope, stated honestly.} This is a null on an enumerated space of $110$
parametrised rules at $n=3,4$, not a proof that no local rule exists. It does
not refute Side~B --- Lenart--Sottile's construction is certainly outside this
space (that is the point of their paper) --- and it is not evidence about their
weight's root-multilinearity. What it does establish is that the obvious
guesses fail, which is why the ``transformation'' was never one line, and it is
the concrete support for Remark~\ref{rem:pred}.

\paragraph{G3.} Verification is at $n\le5$; Theorem~\ref{thm:samuel},
Theorem~\ref{thm:forced} and Proposition~\ref{prop:dichotomy} are proved for
all $n$ (indeed in $\sw$), so this is a check, not a dependency.

\paragraph{G4.} Side~A is held at chunk-read, from my own browse log, not from
a fresh fetch of MO 313951. The mathematical content was independently
re-derived and verified here, so an error in my record would have surfaced as a
failed test; but the \emph{attribution} --- that this is what Samuel wrote ---
rests on that log.

\section{Summary of the answer to Q345}

\begin{enumerate}
\item \textbf{The order is Bruhat order}, graded by Coxeter length. The
absolute-order reading is refuted by $u=321$, $w=231$ in $S_{3}$. So both sides
of Stanley's comparison are Bruhat-chain formulas, and the question was never
about the order.
\item \textbf{Samuel's formula is the Monk-labelled Bruhat chain enumerator.}
In the simple-root basis the coefficient of
$\alpha_{p_{1}}\otimes\dots\otimes\alpha_{p_{m}}$ in $T_{w/u}$ is the number of
saturated chains $u\to w$ whose $i$-th reflection $t_{a_{i}b_{i}}$ satisfies
$a_{i}\le p_{i}<b_{i}$. The tensor algebra is dual-basis bookkeeping; the
modulus $(a,b)+(b,c)=(a,c)$ is the additivity of Monk-label intervals.
\item \textbf{It is a theorem, with a three-line proof}: iterated Monk, then
commutativity of $H^{*}$. Verified exactly at $n=3,4,5$ in all degrees against
an independent Schubert-polynomial instrument, with a negative control.
\item \textbf{The relation is forced}, already in degree $1$, by Monk's rule;
$V$ is the largest quotient of the free space on transpositions in which the
identity can hold. Smallest witness: $t_{13}$ versus $t_{12}+t_{23}$.
\item \textbf{The transformation question is now one checkable property}:
is Lenart--Sottile's chain weight root-multilinear
(Corollary~\ref{cor:q345})? The weak, non-local sense of ``transformation'' is
automatic and empty (Proposition~\ref{prop:dichotomy}(iii)), so this is the
only reading with content. If yes, their theorem is Samuel's identity
specialised along an explicit functional, and the family of all such
specialisations is equivalent to his. If no, no chain-level transformation
exists in either direction. My recorded prediction is \emph{no}
(Remark~\ref{rem:pred}); deciding it needs one deep read, not more
computation.
\end{enumerate}

\noindent In Pak's terms (Q342): this pair is not yet confirmed-independent nor
confirmed-equivalent, but it is no longer vague. The ``is there a
transformation?'' question of 2018 has been replaced by a yes/no property of a
single definition, and one of the two sides is now completely understood.

\end{document}
